% ---
% Reemplaza líneas 2924–2940 de Docs/Technical_Report.tex.
% Decisiones que lo respaldan: H-001, H-010, D-012.
% Capítulo 4 desde el 2026-10-14 (antes §5).
% ---

\newpage
\section{Análisis Exploratorio de Datos}
\label{sec:eda}

% fuente: H-001, H-010, D-012
El presente capítulo aborda la caracterización cuantitativa y técnica
del conjunto de datos CBIS-DDSM (\textit{Curated Breast Imaging Subset of
DDSM}), procesando de manera unificada sus cuatro archivos de metadatos
clínicos y enlazándolos con los 3\,568 estudios mamográficos en formato
DICOM disponibles en disco (recuperados íntegramente tras solventar la
omisión de 84 series mediante la API REST de TCIA).

Lejos de constituir un ejercicio meramente descriptivo, este análisis
exploratorio actúa como el puente empírico fundamental que justifica
y condiciona de forma directa las decisiones de diseño del pipeline
experimental desarrollado en el capítulo~\ref{sec:metodologia}. En
particular, la exploración busca dar respuesta a tres interrogantes
estructurales:
\begin{enumerate}
  \item \textbf{Estructura y balance de clases:} ¿Cuál es el grado de
        asimetría diagnóstica entre benignidad y malignidad en cada
        subconjunto, y qué esquemas de ponderación o métricas de
        evaluación exige dicha distribución?
  \item \textbf{Relevancia de los metadatos clínicos frente a la imagen:}
        ¿Aportan las variables categóricas y ordinales registradas
        (densidad mamaria, sutileza, lateralidad o evaluación BI-RADS)
        señal predictiva utilizable, o debe extraerse la información
        discriminativa directamente del tejido tisular en la imagen?
  \item \textbf{Factibilidad y restricciones de la representación cuántica:}
        ¿Qué características técnicas (resolución espacial, rango dinámico,
        profundidad de bits, alineación de máscaras de contorno y
        ausencia de espaciado físico de píxel) delimitan la viabilidad
        de codificar la información en estados cuánticos de pocos qubits?
\end{enumerate}

El capítulo se estructura en cuatro bloques temáticos: saneamiento y
validación del dataset, análisis de la variable patológica objetivo,
evaluación de correlaciones en variables clínicas, e inspección técnica
de las imágenes DICOM. Finalmente, se sintetizan las implicaciones
directas que estos hallazgos imponen sobre cada módulo de la metodología.
Todos los procedimientos y figuras analíticas descritos en este capítulo
fueron ejecutados en el cuaderno \texttt{1\_EDA.ipynb} del repositorio.
